\chapter{Daily Life with the Guix CLI}
\label{chap:cli-daily-life}

\epigraph{``There are two types of sysadmins: those who have broken production with an update, and those who are about to.''}{--- Murphy's Law of Operations}

\noindent
Now that you understand that \texttt{/gnu/store} is a cryptographic citadel of immutable purity, let us roll up our sleeves and interact with the Guix command-line interface.

Unlike traditional package managers that demand \texttt{sudo} permissions to modify the global system state, Guix package management is entirely \textbf{unprivileged}. Every user on the system can install, upgrade, and remove software in their own user profile without ever asking the sysadmin for permission.

\section{The Non-Declarative CLI: A Foreign Distro Lifeline}

Before typing our first installation command, we must address an essential philosophical question: \emph{How does the Guix command-line interface fit into the functional, declarative vision of GNU Guix?}

The commands introduced in this chapter---\guixcmd{guix install}, \guixcmd{guix remove}, and \guixcmd{guix upgrade}---represent the \textbf{imperative}, or \textbf{non-declarative}, style of package management. In an imperative model, you type sequential commands to mutate your active user profile step by step over time.

It is crucial to understand that \textbf{this non-declarative way of installing packages mainly applies when Guix is used as an auxiliary package manager on a distribution other than Guix System} (what the Guix community terms a \textbf{foreign distribution}, such as Debian, Ubuntu, Fedora, openSUSE, or Arch Linux).

\subsection{Why the Non-Declarative Model Thrives on Foreign Distros}

When you install GNU Guix on top of Ubuntu or Debian, you do not control the base operating system with Guix; your host kernel, systemd init system, display manager, and base utilities are all governed by the host distribution. In this environment, your goals are practical and immediate:
\begin{itemize}
  \item You want cutting-edge software or specific development tools without waiting for your host distribution's slow release cycle.
  \item You lack \texttt{sudo} or root privileges on a shared university cluster or corporate workstation.
  \item You need to run conflicting libraries side-by-side without contaminating the host system's \texttt{/usr/lib}.
\end{itemize}

In this foreign distribution setting, running \guixcmd{guix install} is wonderfully liberating. It acts as a direct, unprivileged replacement for traditional package managers like \texttt{apt} or \texttt{brew}, yet bestows the mathematical superpowers of the Store: atomic transactions, zero dependency conflicts, and instantaneous rollbacks.

\subsection{Why Guix System Purists Avoid \texttt{guix install}}

Conversely, if you are running \textbf{Guix System}---where Guix is the standalone operating system---this non-declarative approach is largely considered an \textbf{anti-pattern} and is discouraged for permanent workstation setups.

Why? Because Guix System is built on the core tenet that \textbf{the file is the machine}:
\begin{itemize}
  \item System packages belong in the \texttt{packages} field of your declarative \texttt{/etc/config.scm} (Chapter~\ref{chap:guix-system}).
  \item User tools and dotfiles belong in a declarative home environment managed by \textbf{Guix Home} (Chapter~\ref{chap:guix-home}) or project-specific manifests (Chapter~\ref{chap:manifests}).
  \item Temporary, disposable tools are best invoked on-demand using ephemeral \guixcmd{guix shell} environments (Chapter~\ref{chap:guix-shell}).
\end{itemize}

If you run \guixcmd{guix install} on Guix System, you create an unversioned, mutable profile under \texttt{\textasciitilde/.guix-profile}. If your hard drive fails or you replicate your environment on another computer, those imperatively installed packages will be completely absent because they were never committed to your declarative Scheme configuration!

\begin{alchemybox}[Non-Declarative CLI vs. Declarative Configuration]
\begin{itemize}
  \item \textbf{Non-Declarative (\texttt{guix install}, \texttt{guix remove})}:
  Best suited for foreign distros (Debian, Ubuntu, Fedora) and fast ad-hoc experimentation. State is tracked via profile generations in \texttt{/var/guix/profiles}, but cannot be fully reproduced on another machine without replaying your shell history.
  \item \textbf{Declarative (\texttt{config.scm}, Guix Home, Manifests)}:
  The standard way on Guix System. State is declared in Scheme files tracked with Git. You reconfigure your system with \guixcmd{guix system reconfigure} or your user environment with \guixcmd{guix home reconfigure}. Any machine can be cloned with 100\% mathematical precision.
\end{itemize}
\end{alchemybox}

Why, then, do we begin our journey with the non-declarative CLI? Because it is the gentlest and most intuitive bridge from traditional package managers. It allows you to grasp profiles, generations, rollbacks, and garbage collection hands-on before we ascend to the declarative peaks of manifests and Guix System.

\section{Installing and Removing Packages}

The most basic operations will look pleasantly familiar, with a subtle functional twist:

\begin{lstlisting}[style=guixcli]
# Search for software (searches names, synopses, and descriptions)
guix search emacs

# Show detailed information about a package
guix show emacs-ripgrep

# Install a package into your active user profile
guix install htop git ripgrep

# Remove a package from your user profile
guix remove htop

# Upgrade all packages in your profile to their latest versions
guix upgrade
\end{lstlisting}

\begin{footgunbox}[The Imperative Trap on Guix System]
If you are reading this while running Guix System, resist the urge to turn \guixcmd{guix install} into your daily routine! Imperatively installed packages live in an isolated profile outside your system's declarative \texttt{config.scm}. Treat the commands in this chapter as mastering the transactional mechanics of profiles and rollbacks, but save your permanent workstation setup for the declarative configurations in Part~II and Part~III.
\end{footgunbox}

\begin{protipbox}[Fast Searching with Recutils]
When you run \guixcmd{guix search}, Guix evaluates the Scheme package records and outputs results formatted as \texttt{recutils} records. You can pipe the output into standard tools or filter fields:
\begin{lstlisting}[style=guixcli]
guix search emacs | grep -E '^name:'
\end{lstlisting}
Or search strictly by package name without scanning full descriptions:
\begin{lstlisting}[style=guixcli]
guix search '^emacs$'
\end{lstlisting}
\end{protipbox}

\section{The Miraculous Undo Button: Rollbacks}

Suppose you ran \guixcmd{guix upgrade} on a Monday morning. The new version of your favorite software contains a catastrophic regression, or perhaps you accidentally uninstalled your text editor right before a deadline.

On an imperative distribution (Ubuntu, Arch, etc.), you would now be combing through cache archives, manually downloading \texttt{.deb} or \texttt{.pkg.tar.zst} files, and praying you don't trigger dynamic linker errors.

On Guix, you simply invoke the undo button:

\begin{lstlisting}[style=guixcli]
# Instant rollback to the previous generation
guix package --roll-back
\end{lstlisting}

\begin{tearsbox}[The 2-Millisecond Rescue]
When you type \guixcmd{guix package --roll-back}, Guix does not download anything. It does not recompile anything. It merely updates the \texttt{\textasciitilde/.guix-profile} symlink to point back to Generation $N-1$. The entire rollback takes approximately 2 milliseconds. You can do it while running a live presentation and nobody in the audience will notice.
\end{tearsbox}

\section{Inspecting and Navigating Generations}

Every transactional change to your profile creates a permanent historical record called a \textbf{Generation}. You can inspect the entire history of your user environment:

\begin{lstlisting}[style=guixcli]
guix package --list-generations
\end{lstlisting}

This displays an informative timeline:

\begin{lstlisting}[style=guixcli]
Generation 1    Oct 12 2025 14:02:11
 + git          2.41.0          out     /gnu/store/7a3...-git-2.41.0
 + ripgrep      13.0.0          out     /gnu/store/8bc...-ripgrep-13.0.0

Generation 2    Oct 14 2025 09:15:30
 + python       3.10.7          out     /gnu/store/9df...-python-3.10.7

Generation 3 (current)  Oct 15 2025 11:22:04
 - ripgrep      13.0.0          out     /gnu/store/8bc...-ripgrep-13.0.0
 + ripgrep      14.0.1          out     /gnu/store/2ea...-ripgrep-14.0.1
\end{lstlisting}

You can switch to any arbitrary point in your history at will:

\begin{lstlisting}[style=guixcli]
# Travel directly to Generation 1
guix package --switch-generation=1
\end{lstlisting}

\section{Multiple Custom Profiles}

Why cram your video editor, your Haskell compiler, your audio synthesizer, and your scientific Python stack into a single bloated user profile?

In Guix, you can instantiate as many independent profiles as you want by passing the \texttt{-p} or \texttt{--profile} flag:

\begin{lstlisting}[style=guixcli]
# Create a dedicated profile for audio production
guix package -p ~/profiles/audio -i ardour jack2 audacity

# Create a dedicated profile for data science
guix package -p ~/profiles/datascience -i python python-numpy python-scipy

# Activate a custom profile in your current shell
source ~/profiles/audio/etc/profile
\end{lstlisting}

When you source \texttt{etc/profile}, Guix automatically configures your environment variables (\texttt{PATH}, \texttt{LIBRARY\_PATH}, \texttt{GUIX\_PYTHONPATH}, etc.) for that specific set of packages.

\begin{protipbox}[Declarative Profiles with Manifests]
While we create custom profiles imperatively in this chapter using \guixcmd{-i}, Guix also lets you instantiate and update profiles \textbf{declaratively} from a Scheme manifest file:
\begin{lstlisting}[style=guixcli]
guix package -p ~/profiles/datascience -m manifest.scm
\end{lstlisting}
This combines the isolation of separate profiles with the reproducibility of source-controlled configuration, which we explore in Chapter~\ref{chap:manifests}.
\end{protipbox}

\section{Garbage Collection and Disk Space Management}

Because Guix keeps old generations around to make rollbacks instantaneous, your \texttt{/gnu/store} will eventually grow larger as you install and upgrade software over months.

\begin{alchemybox}[How Garbage Collection Works]
Guix uses a tracing garbage collector, similar to programming language runtimes (like Java, Go, or Guile).
\begin{enumerate}
  \item \textbf{GC Roots}: Guix scans known roots (active user profiles, system generations, running processes, and pins in \texttt{/var/guix/gcroots}).
  \item \textbf{Liveness Trace}: Any store item that can be reached via symlinks or references from an active root is marked as \textbf{live}.
  \item \textbf{Reclamation}: Any store item that is completely unreferenced is deleted from disk.
\end{enumerate}
\end{alchemybox}

If you simply run \guixcmd{guix gc}, it will only remove packages that are not referenced by \textbf{any} generation of any profile. To reclaim significant space, you first delete historical generations you no longer need:

\begin{lstlisting}[style=guixcli]
# Delete all profile generations older than 30 days
guix package --delete-generations=30d

# Delete all generations except the current one
guix package --delete-generations

# Run the garbage collector to reclaim unused disk space
guix gc

# Free disk space down to a target size (e.g., free at least 10GB)
guix gc -F 10G
\end{lstlisting}

\begin{footgunbox}[The Premature GC Trap]
If you delete your old generations with \guixcmd{guix package --delete-generations} and then immediately run \guixcmd{guix gc}, your previous generations are gone forever. You will no longer be able to roll back to yesterday's profile! Keep at least a few days of generations around unless you are in desperate need of disk space.
\end{footgunbox}

\section{Substitutes: Why You Do Not Have to Compile Everything}

A common fear among newcomers is: \emph{``If Guix is functional and source-based, am I going to spend the next four days compiling Firefox and Chromium from source?''}

The answer is \textbf{no}. Guix uses \textbf{Substitutes} (pre-built binary caches signed by trusted continuous integration servers like \texttt{ci.guix.gnu.org} and \texttt{bordeaux.guix.gnu.org}).

When Guix determines that a derivation hash \texttt{/gnu/store/abc123...-firefox-115.0.drv} is needed, it checks if a trusted substitute server already built that exact derivation. If it exists, Guix simply downloads the pre-built cryptographic narball (normalized archive) and unpacks it into your store. You get all the speed of a binary package manager with all the mathematical purity of a source-based functional system.

\section{The Road Ahead: Beyond Non-Declarative Mutations}

The non-declarative CLI commands covered in this chapter provide an immediate, dependable upgrade for anyone running Guix on a foreign distribution. You can install, upgrade, and rollback software without root privileges and without fear of library collision.

Yet, imperative management has a fundamental ceiling: your environment remains an artifact of whatever sequence of commands you typed into your terminal over the past six months. In the upcoming chapters, we will transcend imperative mutations:
\begin{itemize}
  \item In Chapter~\ref{chap:channels-and-time-travel}, we pin the package tree to precise Git revisions and travel through time.
  \item In Chapter~\ref{chap:guix-shell} and Chapter~\ref{chap:manifests}, we replace permanent profile installations with ephemeral shells and version-controlled manifests.
  \item In Chapter~\ref{chap:guix-system} and Chapter~\ref{chap:guix-home}, we achieve full declarative mastery over entire operating systems and user homes.
\end{itemize}
